\documentclass[11pt]{article}
\usepackage[margin=2.2cm]{geometry}
\usepackage{booktabs,longtable,array}
\usepackage{amsmath,amssymb}
\usepackage[colorlinks=true,linkcolor=blue,urlcolor=blue]{hyperref}
\setlength{\parskip}{4pt}
\setlength{\parindent}{0pt}
\renewcommand{\arraystretch}{1.15}
\newcommand{\dex}{\ensuremath{\,\mathrm{dex}}}

\title{\bfseries Supporting Information\\[4pt]
\large A midpoint-definition artefact reconciles the standard MWC model\\
with single-cell chemotaxis dose responses}
\author{Huan Li\\ \small Shanghai Liangta Biotech Co., Ltd., Shanghai 201899, P.\ R.\ China}
\date{November 2026 (corrected protocol, release v2.0.0)}

\begin{document}
\maketitle

\noindent This Supporting Information accompanies the \emph{corrected} manuscript
(release v2.0.0). It supersedes the v1.x SI in full. Section~1 states the correction;
Sections~2--6 document the corrected computational pipeline, all of which reruns from the
repository root with plain Python~3 (numpy/scipy/matplotlib) against the CC0 FRET archive in
\texttt{data\_external/moore2024\_fret/}. Numbers below are transcribed from
\texttt{results/corrected\_v2/} (regenerated 2026-11-14).

\tableofcontents
\newpage

\section{The correction in one page}

Release v1.0.x of this repository claimed that no parameter set of the standard
MWC-plus-perfect-adaptation class can jointly reproduce the population amplitude table and the
per-cell midpoint distribution of the Moore et al.\ archive (midpoint error $0.622\dex$ at the
amplitude-optimal fit). That claim was produced by two definition mismatches in our own
pipeline, both in \texttt{code54}:

\begin{center}
\begin{tabular}{p{2.2cm} p{5.6cm} p{5.6cm}}
\toprule
 & Observation side & Model side (\texttt{K12\_model}) \\
\midrule
E1 axis & total ligand $T=B+F$ (script docstring; CSV column \texttt{K12\_total\_axis\_uM})
        & returned the foreground increment $F$ \\
E2 response level & interpolated at the absolute value $da=0.5$
        & relative half-maximum (as stated in the v1 Methods text) \\
\bottomrule
\end{tabular}
\end{center}

Consequences: E1 alone contributes $\approx$0.36\dex{} of the $0.622\dex$ (at $B=100\,\mu$M the
model reports $7.3\,\mu$M against the measured $119.6\,\mu$M, an artificial $1.21\dex$);
E2 censors $45$--$68\%$ of cells where the relative convention censors $2/268$. With both
definitions made consistent, a single parameter set reproduces both statistics
($R^2=0.980$, mean midpoint error $0.023\dex$). The v1 headline claims are retracted:
the Pareto-front exclusion, the $N\approx2$--$4$ versus $12$--$39$ cooperativity contradiction,
and the failure of all eight repair hypotheses. Details and the full claim-by-claim table are in
\texttt{ERRATUM.md} at the repository root.

\section{Corrected pipeline map}

\begin{longtable}{p{5.2cm} p{6.2cm} p{3.6cm}}
\toprule
Script & Computes & Output \\
\midrule
\endhead
\texttt{code/recheck\_chemotaxis.py} & the $2\times2$ protocol crossing (axis $\times$ response
level), closed-form vs \texttt{brentq} check ($1.1\times10^{-12}$), Pareto fronts A--D &
\texttt{recheck\_results.json} \\
\texttt{code/sensitivity\_chemotaxis.py} & threshold-fraction and weak-response-floor sensitivity
(13 configurations) & \texttt{sens.json} \\
\texttt{code/fine\_sweep.py} & fine $\lambda$ grid, shows the $\lambda\approx3.5$ acceptance-region
crossing that the v1 seven-point sweep skipped & console tables \\
\texttt{code/deep\_N\_diversity.py} & cooperativity profile likelihood; per-background midpoint
dispersion; sensitivity--responsiveness covariance & \texttt{deep.json} \\
\texttt{code/deep\_confound.py} & saturation and noise-direction confound checks for the
dispersion--gain correlation & console tables \\
\texttt{code/fig\_v2.py}, \texttt{code/fig2\_v3.py} & manuscript figures 1 and 2 &
\texttt{fig1\_v2\_data.json}, PNGs \\
\bottomrule
\end{longtable}

All seven scripts resolve the repository root from their own location
(\texttt{\_\_file\_\_}) and run unmodified on any machine with the archive present.

\section{The four protocols and their Pareto fronts}

Amplitude $R^2$ and mean midpoint error (dex) at the amplitude-optimal weighting ($\lambda=0$)
and at the midpoint-optimal end of the sweep ($\lambda=30$):

\begin{center}
\begin{tabular}{l l l l l}
\toprule
Protocol & $R^2$ ($\lambda{=}0$) & err ($\lambda{=}0$) & $R^2$ ($\lambda{=}30$) & err ($\lambda{=}30$)\\
\midrule
A: v1 protocol (absolute, $F$ axis) & 0.980 & 0.622 & 0.511 & 0.074 \\
B: axis fixed only              & 0.980 & 0.257 & 0.847 & 0.032 \\
C: response level fixed only    & 0.980 & 0.374 & 0.341 & 0.070 \\
D: both fixed (relative, $B{+}F$) & \textbf{0.980} & \textbf{0.023} & 0.978 & 0.016 \\
\bottomrule
\end{tabular}
\end{center}

Under protocol D every swept weighting $\lambda\in[0,30]$ lies inside the v1 acceptance region
($R^2>0.9$, err $<0.1\dex$). There is no trade-off left to explain.

\section{Weber lines under both conventions}

Per-cell midpoints on the total-ligand axis, fit $K_{1/2}=a+bB$:

\begin{center}
\begin{tabular}{l l l l}
\toprule
Convention & Line & Max deviation & Censoring \\
\midrule
Absolute threshold $da=0.5$ & $K_{1/2}=2.29+1.17\,B$ & $0.10\dex$ & 45--68\% \\
Relative half-maximum & $K_{1/2}=0.91+1.06\,B$ & $0.05\dex$ & $<1\%$ \\
\bottomrule
\end{tabular}
\end{center}

The relative-threshold line is straighter and closer to unit slope, uses $2.7\times$ more cells,
and coincides with the quantity the model actually predicts.

\section{Cooperativity is not identified}

Profile likelihood over fixed $N$ (refitting $K_I$, $K_A$, $a_{\max}$): $R^2$ flattens after
$N\approx15$ ($0.977\to0.980$) while $K_A$ runs to its bound for $N\le15$; the data constrain
$N\gtrsim8$ and leave a strong $N$--$K_A$ ridge. $N=39$ is one point on that ridge, not an
estimate. The full table is in \texttt{results/corrected\_v2/deep.json}.

\section{Residual heterogeneity}

Per-cell midpoint dispersion (relative convention, saturated cells only): sd $0.260\dex$ at
$B=0$ falling to $0.022\dex$ at $B=100\,\mu$M; background-centred $\log K_{1/2}$ co-varies with
responsiveness (Spearman $\rho=+0.205$, $p=6.6\times10^{-7}$, $n=579$). Saturation and
noise-direction artefacts are excluded by the checks in \texttt{deep\_confound.py}; what cannot be
excluded without repeated measurements is a noise contribution to the shallow low-background
dose-responses, so the effect is reported as small.

\section{What survives from v1 unchanged}

The adaptation-relay measurement is independent of the rupture analysis and stands: across 2,217
cells and 77,595 pulse events, the pre-pulse baseline is restored within $2.2\%$ of the median
response amplitude after the 5\,s washout, history-independent over a 90-fold dose range
(\texttt{code/chemotaxis\_adaptation\_audit.py}). The prospective mixed-ligand arm is retained in
the manuscript but is no longer used for model-class discrimination: its frozen tolerance band was
defined for the absolute-threshold statistic, and its third reading came from a control arm marked
``not scored'' in the source CSV. Both qualifications are stated in the main text.

\section{Legacy v1 pipeline}

The v1 scripts (\texttt{code49}--\texttt{code70} series) remain in \texttt{code/} for the
historical record; the v1 manuscript and SI are preserved in
\texttt{manuscript/\_archive\_refuted\_v1/}. They document the analysis \emph{as refuted}. Eight of
the legacy scripts carry hard-coded absolute paths from the original run environment and are not
expected to execute unmodified; the corrected pipeline of Section~2 is the supported entry point.

\end{document}
